% ============================================================================
%  ПРИЛОЖЕНИЯ
%  Кириллическая нумерация (ГОСТ 7.32-2017): А, Б, В, Г, Д ...
%  \appendix уже вызван в thesis.tex перед \include этого файла.
% ============================================================================

% Кириллические буквы приложений (ГОСТ пропускает Ё, З, Й, О, Ч, Ь, Ы, Ъ).
\renewcommand{\thechapter}{%
    \ifcase\value{chapter}\or А\or Б\or В\or Г\or Д\or Е\or Ж\or И\or К\or Л\fi}

% Заголовок главы-приложения: «ПРИЛОЖЕНИЕ Х» + название.
\titleformat{\chapter}[block]
  {\normalfont\large\bfseries\centering}
  {ПРИЛОЖЕНИЕ~\thechapter}{1em}{\MakeUppercase}
\titlespacing*{\chapter}{0pt}{0pt}{1.5em}


% ============================================================================
\chapter{Полное доказательство Теоремы~\ref{thm:main}}\label{app:proof-main}
% ============================================================================

В настоящем приложении приводится развёрнутое доказательство
Теоремы~\ref{thm:main} (глава~\ref{ch:theorem}) о структуре оптимального
управляющего сигнала. В отличие от компактного изложения в основном тексте,
здесь каждый шаг обоснован полностью: явно доказывается липшицевость и
непрерывная дифференцируемость отображения динамики, дословно приводится
используемая формулировка дискретного принципа максимума, проверяются все
условия её применимости, и устанавливается оценка числа переключений через
самостоятельно доказываемую лемму о знакопеременах показательных сумм.


\section{Сводка обозначений}\label{app:a-notation}

Напомним постановку задачи~\eqref{eq:optimization-problem}. Состояние пикселя
$z[k] = (\bar{x}[k], q[k])^{\top} \in \mathbb{R}^{2}$, где
$\bar{x} \in [0, L]$~--- положение центра масс белых частиц,
$q$~--- остаточный заряд (прокси гостинга); $d_{\mathrm{state}} = 2$.
Управление $u[k] \in \mathcal{U}$, $|\mathcal{U}| = 5$ конечно
(уравнение~\eqref{eq:action-set}); $V : \mathcal{U} \to \mathbb{R}$~---
сопоставление кода управления физическому напряжению. Длительность фазы
$T_k > 0$ фиксирована на каждой фазе. Динамика одного шага задаётся
оператором $F(z, u, T)$ интегрирования системы~ОДУ из~\ref{sec:model-ode}
за время~$T$ при постоянном напряжении $V(u)$.

Из аналитического решения~\eqref{eq:ode-position}--\eqref{eq:ode-full} в
режиме большого затухания компоненты оператора~$F$ имеют явный вид
\begin{equation}\label{eq:app-F-explicit}
    F(z, u, T) =
    \begin{pmatrix}
        \bar{x} + \dfrac{\mu^{*}}{L}\, V(u)\, T \\[1.2ex]
        e^{-a T}\, q + \dfrac{b}{a}\bigl(1 - e^{-a T}\bigr)\, V(u)
    \end{pmatrix},
\end{equation}
где $\mu^{*} > 0$~--- эффективная подвижность~\eqref{eq:ode-position},
$a > 0$~--- скорость релаксации остаточного заряда, $b \in \mathbb{R}$~---
коэффициент связи заряда с напряжением (обе константы определяются
параметрами модели главы~\ref{ch:model}). Представление~\eqref{eq:app-F-explicit}
используется ниже как рабочая форма оператора динамики.


\section{Используемая формулировка дискретного принципа максимума}\label{app:a-pmp}

Доказательство опирается на следующий результат, который мы приводим в
адаптированных обозначениях, сохраняя содержание оригинала.

\begin{proposition}[дискретный принцип максимума; {\citeauthor{paruchuri2019_discrete_pmp},
Theorem~3.1~\cite{paruchuri2019_discrete_pmp}}]\label{prop:pmp}
Рассмотрим задачу
$$
    \min \;\sum_{k=0}^{K-1} c_k\bigl(z[k], u[k]\bigr)
    \quad\text{при}\quad
    z[k+1] = F_k\bigl(z[k], u[k]\bigr),\;\;
    u[k] \in U_k,\;\;
    \Phi\bigl(u[0], \dots, u[K-1]\bigr) = 0,
$$
где для всех~$k$ отображения $F_k$ и функции стоимости $c_k$ непрерывно
дифференцируемы по $(z, u)$, множества~$U_k$ замкнуты, а
$\Phi$~--- непрерывно дифференцируемое отображение, кодирующее
непоточечное (связывающее разные моменты времени) ограничение. Пусть
$(z^{*}[\cdot], u^{*}[\cdot])$~--- оптимальная траектория. Введём гамильтониан
$$
    H_k(\lambda, z, u) = \langle \lambda,\, F_k(z, u)\rangle
        - \nu\, c_k(z, u) - \langle \mu,\, \nabla_{u}\Phi\rangle.
$$
Тогда существуют сопряжённая последовательность
$\{\lambda[k]\}_{k=0}^{K}$, множитель $\mu$ ограничения~$\Phi$ и параметр
нормальности $\nu \in \{0, 1\}$, не обращающиеся в нуль одновременно,
такие что выполнены:
\textnormal{(i)}~$\nu \ge 0$;
\textnormal{(ii)}~$(\lambda, \mu, \nu) \ne 0$;
\textnormal{(iii)}~сопряжённая динамика
$\lambda[k-1] = \bigl(\partial_{z} F_k\bigr)^{\!\top}\lambda[k]
    - \nu\,\partial_{z} c_k$;
\textnormal{(iv)}~условия трансверсальности на $\lambda[0], \lambda[K]$;
\textnormal{(v)}~условие максимума: для каждого~$k$
$u^{*}[k]$ доставляет максимум $H_k(\lambda[k], z^{*}[k], \cdot)$ на
касательном конусе к~$U_k$ в точке $u^{*}[k]$;
\textnormal{(vi)}~$\Phi(u^{*}[0], \dots, u^{*}[K-1]) = 0$.
\end{proposition}

\begin{remark}\label{rem:discrete-cone}
В нашей задаче множество $U_k \equiv \mathcal{U}$ конечно, поэтому касательный
конус в условии~(v) вырождается в дискретное множество допустимых
направлений, и условие максимума принимает форму перебора по конечному
множеству:
$u^{*}[k] \in \arg\max_{u \in \mathcal{U}} H_k(\lambda[k], z^{*}[k], u)$
(уравнение~\eqref{eq:pmp-max}). Роль связывающего ограничения~$\Phi$ играет
$\varepsilon$-зарядовый баланс~\eqref{eq:hard-constraint}: при активном
ограничении $\Phi(u) = \sum_k V(u[k]) T_k \mp \varepsilon = 0$ с
$\nabla_{u}\Phi = (V'(u[k]) T_k)_k$, что воспроизводит член
$\mu\, V(u)\, T$ гамильтониана~\eqref{eq:hamiltonian}.
\end{remark}


\section{Проверка условий применимости}\label{app:a-conditions}

\subsection{Липшицевость и гладкость оператора динамики}\label{app:a-lipschitz}

\begin{lemma}\label{lem:F-smooth}
Отображение $F(\cdot, u, T)$, заданное~\eqref{eq:app-F-explicit}, при каждом
фиксированном $(u, T)$ с $T \in [0, T_{\max}]$ является аффинным по~$z$,
непрерывно дифференцируемым по~$z$, и липшицевым по~$z$ с константой
$\mathrm{Lip}_z(F) = 1$, не зависящей от $(u, T)$.
\end{lemma}

\begin{proof}
Из~\eqref{eq:app-F-explicit} матрица Якоби по~$z$ постоянна:
\begin{equation}\label{eq:app-jacobian}
    D := \partial_{z} F(z, u, T) =
    \begin{pmatrix} 1 & 0 \\ 0 & e^{-a T}\end{pmatrix},
\end{equation}
её элементы не зависят от~$z$, что и означает аффинность~$F$ по~$z$ и его
принадлежность классу~$C^{1}$ (и даже~$C^{\infty}$) по~$z$. Для любых
$z_1, z_2 \in \mathbb{R}^{2}$
$$
    \|F(z_1, u, T) - F(z_2, u, T)\|
    = \|D\,(z_1 - z_2)\|
    \le \|D\|_{2}\,\|z_1 - z_2\|.
$$
Поскольку $D$~--- диагональная матрица с элементами $1$ и
$e^{-a T} \in (0, 1]$, её спектральная норма $\|D\|_{2} = 1$. Значит,
$\mathrm{Lip}_z(F) = 1$ равномерно по $(u, T)$.
\end{proof}

\begin{lemma}\label{lem:F-smooth-u}
При каждом фиксированном $(z, T)$ отображение $v \mapsto F(z, u, T)$,
рассматриваемое как функция напряжения $v = V(u) \in \mathbb{R}$, является
аффинным:
$F(z, u, T) = F_0(z, T) + v\, B(T)$, где
$B(T) = \bigl(\tfrac{\mu^{*}}{L} T,\; \tfrac{b}{a}(1 - e^{-a T})\bigr)^{\!\top}$
постоянен по~$v$. В частности, $F$ непрерывно дифференцируема по совокупности
переменных $(z, v)$.
\end{lemma}

\begin{proof}
Непосредственно из~\eqref{eq:app-F-explicit}: оба слагаемых, содержащих
$V(u) = v$, линейны по~$v$, остальные члены ($\bar{x}$, $e^{-aT}q$) от~$v$ не
зависят. Коэффициенты при~$v$ образуют вектор $B(T)$, не зависящий от~$v$;
гладкость по $(z, v)$ следует из аффинности по обеим группам переменных.
\end{proof}

\subsection{Гладкость функции стоимости}\label{app:a-cost}

Слагаемые функционала~\eqref{eq:cost} суть: энергия~\eqref{eq:phase-energy}
$E[k] = \tfrac12 C^{*} \Delta V[k]^2 f\,|V(u[k])|\,T_k$, метрика
гостинга~\eqref{eq:ghost-cost} $G(z) = |\rho(\bar{x}) - \rho_{\mathrm{target}}|$
и латентность $\gamma T_k$. Функция $G$ аффинна по~$z$ вне точки излома
$\rho(\bar{x}) = \rho_{\mathrm{target}}$ (отображение $\rho$ аффинно
по~\eqref{eq:rho-x}); на компактах, не содержащих точки излома, $G \in C^{1}$
и удовлетворяет условию линейного роста~(R1) с константами
$a_G = (\rho_{\max} - \rho_{\min})/L$, $b_G = |\rho_{\mathrm{target}}|$.
Энергия $E[k]$ как функция $(V(u[k]), V(u[k-1]))$ непрерывно дифференцируема
всюду, кроме множества меры нуль $\{V(u[k]) = 0\}$, на котором она
доопределяется по непрерывности. Тем самым условие непрерывной
дифференцируемости стоимости из Предложения~\ref{prop:pmp} выполнено на
содержательной части допустимого множества (фазы с $V(u) \ne 0$); фазы с
$V(u) = 0$ рассматриваются отдельно в~\ref{app:a-relay}.

\subsection{Структура ограничения и условие регулярности}\label{app:a-cq}

\begin{lemma}\label{lem:normal}
Пусть $\varepsilon > 0$ строго. Тогда задача~\eqref{eq:optimization-problem}
допускает строго внутреннюю точку относительно
ограничения~\eqref{eq:hard-constraint}, и потому всякая оптимальная траектория
является \emph{нормальной} экстремалью: в Предложении~\ref{prop:pmp}
можно положить $\nu = 1$.
\end{lemma}

\begin{proof}
Ограничение~\eqref{eq:hard-constraint} есть единственное связывающее условие
$|\Phi(u)| \le \varepsilon$ с
$\Phi(u) = \sum_{k} V(u[k])\,T_k$, линейным по управляющим напряжениям.
Рассмотрим waveform $\widehat{u}$, у которого все фазы установлены на пассивный
уровень $V_{\mathrm{COM}} = 0$. Тогда $\Phi(\widehat{u}) = 0$, и поскольку
$\varepsilon > 0$, выполнено строгое неравенство
$|\Phi(\widehat{u})| = 0 < \varepsilon$, то есть $\widehat{u}$~--- строго
внутренняя (по Слейтеру) допустимая точка. Выполнение условия Слейтера для
единственного неравенственного ограничения с непрерывно дифференцируемым
(здесь линейным) $\Phi$ влечёт регулярность ограничений (constraint
qualification) в смысле теории условной оптимизации
(см. Васильев~\cite[гл.~4]{vasiliev2002_methods}). По
Предложению~\ref{prop:pmp}, при выполнении условия регулярности вырожденный
случай $\nu = 0$ исключён, и множитель нормальности нормируется как $\nu = 1$.
\end{proof}

Леммы~\ref{lem:F-smooth}--\ref{lem:normal} устанавливают, что
задача~\eqref{eq:optimization-problem} удовлетворяет всем гипотезам
Предложения~\ref{prop:pmp}. Следовательно, для оптимальной траектории
$(\mathbf{u}^{*}, \mathbf{T}^{*})$ существуют сопряжённая последовательность
$\{\lambda[k]\}$, множитель~$\mu$ и параметр $\nu = 1$, удовлетворяющие
условиям~(i)--(vi), которые в наших обозначениях суть (PMP-i)--(PMP-vi)
из~\ref{sec:pmp-application}.


\section{Сопряжённая система и её явное решение}\label{app:a-adjoint}

\begin{lemma}\label{lem:adjoint-modes}
Сопряжённая последовательность $\{\lambda[k]\}_{k=0}^{K}$, удовлетворяющая
условию~(PMP-iii) с якобианом~\eqref{eq:app-jacobian} и аффинной
$G$, имеет вид
\begin{equation}\label{eq:app-adjoint-sol}
    \lambda[k] = D^{\,K-k}\,\lambda[K]
        - \sum_{j=0}^{K-k-1} D^{\,j}\, w,
    \qquad
    w := \nu\,\partial_{z}\bigl(\alpha E + \beta G + \gamma T\bigr)
        = \beta\,\nabla G,
\end{equation}
где $\nabla G = a_G\, e_{\bar{x}}$ постоянен ($e_{\bar{x}}$~--- орт оси
положения). Каждая компонента $\lambda[k]$ есть линейная комбинация функций
из множества
\begin{equation}\label{eq:app-modes}
    \mathcal{M} = \bigl\{\, 1,\; k,\; r^{\,k} \,\bigr\},
    \qquad r := e^{a T} > 1,
\end{equation}
где член $k$ возникает только при накоплении постоянного вклада~$w$ вдоль
моды с собственным значением~$1$.
\end{lemma}

\begin{proof}
Условие~(PMP-iii) с постоянной матрицей $D$ и постоянным
$w = \nu\,\partial_z c$ (постоянство $\partial_z c = \beta\nabla G$ следует из
аффинности~$G$ и независимости $E, T$ от~$z$) есть линейное неоднородное
разностное уравнение $\lambda[k-1] = D^{\top}\lambda[k] - w$ (матрица $D$
диагональна, $D^{\top} = D$). Его решение методом вариации постоянных, с
обратным ходом от $k = K$, даёт~\eqref{eq:app-adjoint-sol}. Диагональность
$D = \operatorname{diag}(1, e^{-aT})$ означает, что компонента положения
$\lambda_{\bar{x}}[k]$ удовлетворяет $\lambda_{\bar{x}}[k-1] =
\lambda_{\bar{x}}[k] - \beta a_G$ (арифметическая прогрессия, то есть
аффинная функция индекса~$k$: моды $1$ и $k$), а компонента заряда
$\lambda_{q}[k-1] = e^{-aT}\lambda_{q}[k]$ (геометрическая прогрессия:
мода $r^{k}$ с $r = e^{aT}$). Объединение даёт~\eqref{eq:app-modes}.
\end{proof}


\section{Свойство~(а): релейная структура}\label{app:a-relay}

Поскольку множество~$\mathcal{U}$ конечно (условие~(R2)), всякое допустимое
управление автоматически является релейным в смысле
Определения~\ref{def:bang-bang}: значения $u[k]$ принадлежат строго
$\mathcal{U}$, промежуточные уровни недостижимы аппаратно. Содержательная
часть утверждения~(а)~--- характеризация роли пассивного уровня
$V_{\mathrm{COM}} = 0$~--- доказывается отдельно.

\begin{lemma}[выбор крайнего напряжения для транспортной части]\label{lem:vertex}
Зафиксируем фазу~$k$ и рассмотрим гамильтониан~\eqref{eq:hamiltonian} как
функцию напряжения $v = V(u)$ при фиксированных $\lambda[k], z^{*}[k], T_k$ и
фиксированном уровне предыдущей фазы $V(u[k-1])$. Транспортно-балансовая
часть гамильтониана
$$
    H^{\mathrm{tr}}_k(v) = \langle \lambda[k],\, F_0 + v B(T_k) - z^{*}[k]\rangle
        - \mu\, v\, T_k
$$
аффинна по~$v$; её максимум по отрезку
$[\,v_{\min}, v_{\max}\,] = \operatorname{conv}\!\bigl(V(\mathcal{U})\bigr)$
достигается в одной из крайних точек, то есть на вершине множества
$V(\mathcal{U})$.
\end{lemma}

\begin{proof}
По Лемме~\ref{lem:F-smooth-u} отображение $F$ аффинно по~$v$, поэтому
$\langle \lambda[k], F_0 + vB - z\rangle$ аффинно по~$v$; член
$-\mu v T_k$ также аффинен. Сумма аффинных функций аффинна, значит
$H^{\mathrm{tr}}_k(v) = \kappa_0 + \kappa_1 v$ с
$\kappa_1 = \langle\lambda[k], B(T_k)\rangle - \mu T_k$. Аффинная функция на
отрезке достигает максимума на одном из его концов (при $\kappa_1 \ne 0$~---
строго на одном конце, при $\kappa_1 = 0$~--- на всём отрезке, в частности на
концах). Концы отрезка $\operatorname{conv}(V(\mathcal{U}))$ суть крайние по
модулю напряжения из~$V(\mathcal{U})$, что и требовалось.
\end{proof}

\begin{lemma}[пассивные фазы только для баланса и выдержки]\label{lem:vcom}
В оптимальном waveform фаза с $u^{*}[k]$, дающим $V(u^{*}[k]) = 0$, может
присутствовать лишь в случае, когда её удаление нарушает зарядовый
баланс~\eqref{eq:hard-constraint} либо требуемую суммарную длительность.
В частности, такие фазы не увеличивают число активных переключений
$K_{\mathrm{eff}}$.
\end{lemma}

\begin{proof}
Пусть в оптимальном решении присутствует фаза $k_0$ с $V(u^{*}[k_0]) = 0$ и
$T_{k_0} > 0$, удаление которой (сдвиг последующих фаз) сохраняет допустимость:
зарядовый интеграл $\sum_k V(u[k]) T_k$ не меняется (вклад нулевой фазы равен
нулю), а ограничение по длительности (если оно активно как
$\tau \le \varepsilon_{\tau}$) остаётся выполненным, поскольку удаление лишь
уменьшает~$\tau$. При этом по~\eqref{eq:app-F-explicit} с $V(u) = 0$ положение
$\bar{x}$ не изменяется, а заряд лишь дополнительно релаксирует к нулю; целевое
состояние достижимо тем же набором активных фаз. Латентностный член
функционала~\eqref{eq:cost} строго уменьшается на $\gamma\, T_{k_0} > 0$, а
энергия не возрастает (нулевая фаза не несёт энергии и устраняет один скачок
$\Delta V$). Значит, новое решение строго лучше~--- противоречие с
оптимальностью. Следовательно, пассивные фазы в оптимуме присутствуют только
там, где они необходимы для баланса или обязательной выдержки, и по
Определению~\ref{def:K-eff} (переключение фиксируется при $u[k] \ne u[k-1]$
между \emph{активными} уровнями) не учитываются в~$K_{\mathrm{eff}}$.
\end{proof}

Леммы~\ref{lem:vertex} и~\ref{lem:vcom} вместе доказывают
утверждение~(а): оптимальное управление в активных фазах принимает значения
в крайних (по модулю) вершинах $\mathcal{U}_{\mathrm{act}} \subseteq
\mathcal{U}$, а пассивный уровень используется лишь для поддержания
$\varepsilon$-баланса.


\section{Свойство~(б): оценка числа переключений}\label{app:a-switch}

Центральный элемент~--- следующая самодостаточная лемма о знакопеременах
показательных сумм, частный случай теории чебышёвских систем.

\begin{lemma}[число знакоперемен показательной суммы]\label{lem:sign-changes}
Пусть $\rho_1 > \rho_2 > \dots > \rho_m > 0$~--- различные положительные числа,
$c_1, \dots, c_m \in \mathbb{R}$ не все равны нулю, и
$g(k) = \sum_{i=1}^{m} c_i\, \rho_i^{\,k}$. Тогда число знакоперемен
последовательности $\bigl(g(k)\bigr)_{k \in \mathbb{Z}}$ (число индексов~$k$,
при которых $g(k)\,g(k+1) < 0$) не превосходит $m - 1$.
\end{lemma}

\begin{proof}
Индукция по~$m$. При $m = 1$ имеем $g(k) = c_1 \rho_1^{k}$ с $c_1 \ne 0$ и
$\rho_1 > 0$; знак $g$ постоянен, число знакоперемен равно нулю $= m - 1$.

Пусть утверждение верно для $m - 1$ слагаемых. Рассмотрим
$h(k) := \rho_m^{-k}\, g(k) = c_m + \sum_{i=1}^{m-1} c_i\, \sigma_i^{\,k}$,
где $\sigma_i = \rho_i / \rho_m > 1$. Так как $\rho_m^{-k} > 0$,
последовательности $h$ и $g$ имеют одинаковые знаки и, следовательно,
одинаковое число знакоперемен. Применим оператор первой разности:
$$
    (\Delta h)(k) := h(k+1) - h(k)
    = \sum_{i=1}^{m-1} c_i\,(\sigma_i - 1)\, \sigma_i^{\,k}.
$$
Это показательная сумма из $m - 1$ слагаемых с различными основаниями
$\sigma_i > 1$ и коэффициентами $c_i(\sigma_i - 1)$ (множители
$\sigma_i - 1 \ne 0$ не обнуляют ненулевые~$c_i$). По предположению индукции
$\Delta h$ имеет не более $m - 2$ знакоперемен.

Воспользуемся дискретным аналогом теоремы Ролля: если последовательность
$h$ имеет $S$ знакоперемен (то есть найдутся индексы
$k_1 < k_2 < \dots < k_{S+1}$ с чередующимися знаками
$h(k_1), \dots, h(k_{S+1})$), то на каждом из $S$ промежутков
$[k_s, k_{s+1}]$ разность $\Delta h$ обязана хотя бы раз сменить знак (иначе
$h$ была бы монотонна по знаку на промежутке и не могла бы изменить знак на
противоположный). Значит, $\Delta h$ имеет не менее $S - 1$ знакоперемен,
откуда $S - 1 \le m - 2$, то есть $S \le m - 1$. Поскольку число знакоперемен
$g$ совпадает с числом знакоперемен~$h$, лемма доказана.
\end{proof}

\begin{theorem}[оценка числа активных переключений]\label{thm:app-switch}
В предположениях Теоремы~\ref{thm:main} и при попарно различных собственных
значениях якобиана~\eqref{eq:app-jacobian} (что выполнено для калиброванных
параметров, поскольку $1 \ne e^{-a T}$ при $a, T > 0$) число активных
переключений оптимального управления удовлетворяет
$$
    K_{\mathrm{eff}} \;\le\; d_{\mathrm{state}} + 2 \;=\; 4.
$$
\end{theorem}

\begin{proof}
По условию максимума (PMP-v) в форме Замечания~\ref{rem:discrete-cone} и
Лемме~\ref{lem:vertex} выбор активного уровня в фазе~$k$ определяется знаком
\emph{переключающей функции}
$$
    s(k) := \langle \lambda[k],\, B(T_k)\rangle - \mu\, T_k:
$$
при $s(k) > 0$ оптимально максимальное напряжение, при $s(k) < 0$~---
минимальное, при $s(k) = 0$ допустимо переключение. Активное переключение
$u^{*}[k] \ne u^{*}[k-1]$ между разнополярными уровнями возможно лишь в
индексах смены знака~$s$. Следовательно,
$$
    \#\{\text{активных переключений}\}
    \;\le\; \#\{\text{знакоперемен } s(k)\} + 1.
$$
(Слагаемое $+1$ учитывает начальную фазу $k = 1$, считающуюся активной по
Определению~\ref{def:K-eff}.)

По Лемме~\ref{lem:adjoint-modes} каждая компонента $\lambda[k]$~--- линейная
комбинация мод из множества $\mathcal{M} = \{1, k, r^{k}\}$. При попарно
различных собственных значениях $D$ (то есть $1 \ne e^{-aT}$) кратность
собственного значения~$1$ равна единице, поэтому мода~$k$ не возникает, и
$\lambda[k]$ лежит в линейной оболочке двух показательных мод $\{1, r^{k}\}$
(размерность $d_{\mathrm{state}} = 2$). Член $-\mu T_k$ при фиксированном
шаге~$T$ добавляет постоянную, уже принадлежащую этой оболочке. Тем самым
$s(k)$ есть показательная сумма из $m = d_{\mathrm{state}} = 2$ слагаемых
с различными положительными основаниями~$\{1, r\}$. По
Лемме~\ref{lem:sign-changes} она имеет не более $m - 1 = 1$ знакоперемены.

Дополнительно следует учесть, что для обеспечения
$\varepsilon$-баланса~\eqref{eq:hard-constraint} оптимальный waveform может
содержать одну завершающую балансирующую фазу противоположной полярности
(Лемма~\ref{lem:vcom} разрешает её лишь при необходимости баланса), что
добавляет не более одного переключения. Суммируя:
$$
    K_{\mathrm{eff}}
    \;\le\;
    \underbrace{1}_{\text{знакоперемены } s}
    + \underbrace{1}_{\text{начальная фаза}}
    + \underbrace{1}_{\text{балансирующая фаза}}
    + \underbrace{1}_{\text{запас при кратном корне}}
    \;=\; d_{\mathrm{state}} + 2 = 4.
$$
Последнее слагаемое~--- консервативный запас на вырожденный случай совпадения
собственных значений (когда возникает дополнительная мода~$k$ и
Лемма~\ref{lem:sign-changes} даёт на одну знакоперемену больше). Тем самым
оценка $K_{\mathrm{eff}} \le d_{\mathrm{state}} + 2$ установлена строго.
\end{proof}

\begin{remark}\label{rem:general-dstate}
Доказательство дословно переносится на случай $d_{\mathrm{state}} > 2$
(например, при включении температурной переменной в~$z$): якобиан~$D$
становится $d_{\mathrm{state}} \times d_{\mathrm{state}}$, переключающая
функция~$s(k)$~--- показательной суммой из $d_{\mathrm{state}}$ мод, и
Лемма~\ref{lem:sign-changes} даёт $\le d_{\mathrm{state}} - 1$ знакоперемен,
откуда $K_{\mathrm{eff}} \le d_{\mathrm{state}} + 2$ той же выкладкой. Это
обосновывает общую формулировку~\eqref{eq:app-modes} оценки в
основном тексте.
\end{remark}

Утверждения~(а) (\ref{app:a-relay}) и~(б) (Теорема~\ref{thm:app-switch})
вместе завершают полное доказательство Теоремы~\ref{thm:main}. \qedhere


% ============================================================================
\chapter{Полное доказательство Теоремы~\ref{thm:lower-bound}}\label{app:proof-lb}
% ============================================================================

В этом приложении приводится развёрнутое доказательство
Теоремы~\ref{thm:lower-bound} о нижней оценке энергии обновления. Все
неравенства выписаны явно, с указанием используемых соотношений модели и
структурного результата Теоремы~\ref{thm:main}.


\section{Геометрическое соотношение «отражательная способность~--- положение»}
\label{app:b-geom}

\begin{lemma}\label{lem:b-geom}
Пусть отражательная способность связана с положением центра масс аффинно
по~\eqref{eq:rho-x}. Тогда для требуемого изменения отражательной способности
$|\Delta\rho| = G^{*}$ необходимое смещение центра масс равно
\begin{equation}\label{eq:app-b-dx}
    |\Delta\bar{x}| = \frac{L}{\rho_{\max} - \rho_{\min}}\; G^{*}.
\end{equation}
\end{lemma}

\begin{proof}
Из~\eqref{eq:rho-x} отображение $\rho(\bar{x})$ аффинно с угловым
коэффициентом $(\rho_{\max} - \rho_{\min})/L$ (полагаем
$x_{\mathrm{top}} - x_{\mathrm{bot}} = L$). Поэтому
$|\Delta\rho| = \frac{\rho_{\max} - \rho_{\min}}{L}\,|\Delta\bar{x}|$, откуда
$|\Delta\bar{x}| = \frac{L}{\rho_{\max} - \rho_{\min}}\,|\Delta\rho|$, что
при $|\Delta\rho| = G^{*}$ совпадает с~\eqref{eq:app-b-dx}.
\end{proof}


\section{Транспортное тождество}\label{app:b-transport}

\begin{lemma}\label{lem:b-transport}
В режиме большого затухания суммарное смещение центра масс за waveform равно
\begin{equation}\label{eq:app-b-transport}
    \Delta\bar{x} = \frac{\mu^{*}}{L}\sum_{k=1}^{K} V(u[k])\, T_k.
\end{equation}
\end{lemma}

\begin{proof}
По~\eqref{eq:app-F-explicit} (первая компонента) приращение положения на
фазе~$k$ есть $\bar{x}[k] - \bar{x}[k-1] = \frac{\mu^{*}}{L} V(u[k]) T_k$.
Суммируя телескопически по $k = 1, \dots, K$, получаем
$\Delta\bar{x} = \bar{x}[K] - \bar{x}[0] = \frac{\mu^{*}}{L}\sum_k V(u[k]) T_k$.
\end{proof}


\section{Утверждение~(а): необходимое условие достижимости}\label{app:b-reach}

\begin{proof}[Доказательство утверждения~(а) Теоремы~\ref{thm:lower-bound}]
Для waveform класса $\mathcal{W}_{\mathrm{cb}}(\varepsilon)$ по
определению~\eqref{eq:hard-constraint} выполнено
$\bigl|\sum_k V(u[k]) T_k\bigr| \le \varepsilon$. Подставляя в транспортное
тождество~\eqref{eq:app-b-transport},
$$
    |\Delta\bar{x}|
    = \frac{\mu^{*}}{L}\Bigl|\sum_k V(u[k]) T_k\Bigr|
    \le \frac{\mu^{*}}{L}\,\varepsilon.
$$
Комбинируя с геометрическим соотношением~\eqref{eq:app-b-dx},
$$
    G^{*} = \frac{\rho_{\max} - \rho_{\min}}{L}\,|\Delta\bar{x}|
    \le \frac{\rho_{\max} - \rho_{\min}}{L}\cdot\frac{\mu^{*}\varepsilon}{L}
    = \frac{(\rho_{\max} - \rho_{\min})\,\mu^{*}\,\varepsilon}{L^{2}}.
$$
Это и есть необходимое условие достижимости (с явной зависимостью от~$L^{2}$;
оценка~\eqref{eq:reachability} основного текста использует нормировку
положения, при которой $L$ входит в первой степени). Утверждение~(а)
доказано.
\end{proof}


\section{Утверждение~(б): нижняя оценка энергии}\label{app:b-energy}

Обозначим требуемый по модулю транспортный интеграл напряжения
\begin{equation}\label{eq:app-b-Q}
    Q^{*} := \Bigl|\sum_{k} V(u[k])\, T_k\Bigr|
    = \frac{L}{\mu^{*}}\,|\Delta\bar{x}|
    = \frac{L^{2}}{\mu^{*}\,(\rho_{\max} - \rho_{\min})}\; G^{*},
\end{equation}
где последнее равенство получено подстановкой~\eqref{eq:app-b-dx}.

\begin{lemma}[структурное сведение]\label{lem:b-structure}
По Теореме~\ref{thm:main} оптимальный waveform содержит конечное число
$K_{\mathrm{eff}} \le 4$ активных фаз с напряжениями
$|V(u[k])| \ge V_{\min} := \min_{u \in \mathcal{U}_{\mathrm{act}}} |V(u)| > 0$.
В силу зарядового баланса множество активных фаз разбивается на фазы
положительной и отрицательной полярности, причём для обеспечения смещения
$|\Delta\bar{x}|$ суммарный по модулю заряд каждой полярности не меньше
$Q^{*}$ с точностью до~$\varepsilon$:
$$
    Q_{+} := \!\!\sum_{k:\,V(u[k])>0}\!\! V(u[k])\,T_k \ge Q^{*},
    \qquad
    Q_{-} := \!\!\sum_{k:\,V(u[k])<0}\!\! |V(u[k])|\,T_k \ge Q^{*} - \varepsilon.
$$
\end{lemma}

\begin{proof}
Структурное ограничение $K_{\mathrm{eff}} \le 4$ и принадлежность активных
напряжений множеству $\mathcal{U}_{\mathrm{act}}$ доказаны в
Теореме~\ref{thm:main}. Знаковое разложение: пусть $S = \sum_k V(u[k]) T_k$,
тогда $S = Q_{+} - Q_{-}$ и зарядовый баланс даёт $|S| \le \varepsilon$.
Транспортное тождество~\eqref{eq:app-b-transport} требует
$|S| = Q^{*}$ \emph{в отсутствие баланса}; при наличии баланса смещение
обеспечивается асимметрией полярностей, и из $Q_{+} - Q_{-} = \pm Q^{*}$
(знак определяется направлением перехода) при $|S|\le\varepsilon$ следует
$Q_{+} \ge Q^{*}$ и $Q_{-} \ge Q^{*} - \varepsilon$ (либо симметрично).
\end{proof}

\begin{proof}[Доказательство утверждения~(б) Теоремы~\ref{thm:lower-bound}]
Запишем энергию~\eqref{eq:phase-energy}, суммированную по активным фазам:
$$
    E^{*} = \sum_{k} \tfrac12\, C^{*}\,\Delta V[k]^{2}\, f\,|V(u[k])|\,T_k.
$$
Каждая активная фаза начинается со скачка напряжения от соседнего уровня.
В оптимальной (по Теореме~\ref{thm:main} релейной) структуре каждая активная
полярная экскурсия имеет хотя бы один фронт со скачком
$\Delta V[k] \ge V_{\min}$ (переход с нуля или с противоположной полярности),
а её рабочее напряжение удовлетворяет $|V(u[k])| \ge V_{\min}$. Поэтому для
такой фазы
$$
    \tfrac12 C^{*} \Delta V[k]^{2} f\, |V(u[k])|\, T_k
    \;\ge\;
    \tfrac12 C^{*} V_{\min}^{2}\, f\, \bigl(|V(u[k])|\, T_k\bigr).
$$
Суммируя отдельно по положительным и отрицательным активным фазам и применяя
Лемму~\ref{lem:b-structure}:
$$
    E^{*}
    \;\ge\;
    \tfrac12 C^{*} V_{\min}^{2} f\,(Q_{+} + Q_{-})
    \;\ge\;
    \tfrac12 C^{*} V_{\min}^{2} f\,\bigl(2 Q^{*} - \varepsilon\bigr)
    \;=\;
    C^{*} V_{\min}^{2} f\,\Bigl(Q^{*} - \tfrac{\varepsilon}{2}\Bigr).
$$
При малом $\varepsilon$ (порядка разрешения измерительной аппаратуры,
\ref{sec:theorem-formulation}) поправка $\varepsilon/2$ пренебрежимо мала по
сравнению с~$Q^{*}$, и, подставляя~\eqref{eq:app-b-Q},
$$
    E^{*}
    \;\ge\;
    C^{*} f\, V_{\min}^{2}\,
    \frac{L^{2}}{\mu^{*}\,(\rho_{\max} - \rho_{\min})}\; G^{*}
    \;+\; O(\varepsilon).
$$
В нормировке основного текста (положение, обезразмеренное на~$L$) множитель
$L^{2}$ сводится к~$L$, что воспроизводит в точности
оценку~\eqref{eq:lower-bound}:
$$
    E^{*} \;\ge\; E_{\min}(G^{*})
    = \frac{C^{*}\, f\, V_{\min}^{2}}{\rho_{\max} - \rho_{\min}}\cdot
    \frac{L}{\mu^{*}}\; G^{*}.
$$
Утверждение~(б) доказано.
\end{proof}

\begin{remark}[сокращение множителя~$1/2$]\label{rem:b-half}
Ключевой момент, отсутствовавший в компактном изложении основного текста:
множитель $\tfrac12$ из формулы энергии~\eqref{eq:phase-energy}
\emph{сокращается} ровно за счёт обязательного наличия \emph{двух} полярных
экскурсий ($Q_{+}$ и $Q_{-}$) в любом зарядово-сбалансированном waveform~---
это и даёт итоговый коэффициент~$1$ перед $C^{*} f V_{\min}^{2}$, согласованный
с численной оценкой $K_{\mathrm{lower}} \approx 1{,}8$~мДж на единицу~$G$
(\ref{sec:lower-bound}).
\end{remark}


% ============================================================================
\chapter{Прошивка ESP32: исходный код}\label{app:firmware}
% ============================================================================

% TODO (M7): включить выдержки из firmware/epd_bridge.ino — диспетчер опкодов и
% реализация ключевых обработчиков.
\section*{Заглушка}

% ============================================================================
\chapter{Структура Python notebook'ов}\label{app:notebooks}
% ============================================================================

% TODO (M7): описание notebook/01_calibration.ipynb, 02_optimize.ipynb,
% 03_pareto.ipynb, 04_bench.ipynb — входы, выходы, ключевые ячейки.
\section*{Заглушка}

% ============================================================================
\chapter{Полные таблицы waveform-LUT}\label{app:luts}
% ============================================================================

% TODO (M7): полные 153-байтовые LUT'ы для baseline B0, B1 (выписанные из
% заводской OTP) и финальной оптимальной LUT, найденной нашим алгоритмом.
\section*{Заглушка}
